\Qhead{6.2}{Априорные сведения о сигнале и шуме. Наблюдение и сообщение}{Априорные сведения о сигнале и шуме. Наблюдение и сообщение. Задачи интерполяции, фильтрации и экстраполяции.}
\Subfile{6.2.1}{sections/06_detection/q02/s01}

\recap{Задачу определяют неизвестное сообщение, доступные данные и критерий ошибки. Причинная фильтрация, прогноз и сглаживание различаются временем данных. При квадратичной цене оптимально апостериорное среднее.}
\board{Ось времени с тремя задачами; разложение условной квадратичной ошибки.}
\oral{Почему оценка прошлого после получения будущих данных может быть точнее оценки в реальном времени?}
\bait{Можно оставить различие среднего и моды: максимум плотности не обязан минимизировать среднеквадратичную ошибку.}
\sources{Сёрккя и Солин~\cite{sarkka}, глава~10, постановки фильтрации и сглаживания; Биленко и соавт.~\cite{bilenko}, §5.8.}
